\documentclass[UTF8]{ctexart}

\usepackage{amscd,amsthm,amsfonts,amssymb,amsmath}
\usepackage[colorlinks=true]{hyperref}
\usepackage[all]{xy}
\usepackage{mathrsfs}
\usepackage{tikz}
\usepackage{bm}
\usepackage{geometry}
\usepackage{xcolor}
\usepackage{eso-pic}
% \usepackage{CJK}
\usepackage{arydshln}
% \usepackage[11pt]{moresize}



\newtheorem{thm}{定理}
\newtheorem{cor}[thm]{推论}
\newtheorem{lem}[thm]{引理}
\newtheorem{prop}[thm]{命题}
\newtheorem{defn}[thm]{定义}
\newtheorem{ques}[thm]{问题}
\newtheorem{eg}[thm]{例题}
\newtheorem{ex}[thm]{习题}
\newtheorem{rmk}[thm]{注释}
\newtheorem{ntn}[thm]{记号}




\begin{document}


\title{习题课材料（七）}
\date{}

\maketitle


%\noindent{\Large\textbf{注: 带 $\heartsuit $号的习题有一定的难度、比较耗时, 请量力为之.}}

\begin{ex}
  奇数阶反对称矩阵不可逆.
\end{ex}


\begin{ex}
  对分块上三角矩阵 \(X=\begin{bmatrix}A & C \\ O & B\end{bmatrix}\),
  证明 \(\det X = \det A \cdot \det B\).
\end{ex}


\begin{ex}
  设 \(A\), \(B\) 分别为 \(m\times n\), \(n\times m\) 矩阵.
  证明 \(\det(I_m + AB) = \det(I_n+BA)\).
\end{ex}

 

\begin{ex}[\(\heartsuit\)]
  设 \(T = [\boldsymbol{t}_1,\ldots,\boldsymbol{t}_n]\) 为 \(n\) 阶方阵.
  证明 Hadamard 不等式
  \[
    |\det T| \leq \|\boldsymbol{t}_1\| \cdots \|\boldsymbol{t}_n\|.
  \]
\end{ex}




\begin{ex}
  计算行列式
  \begin{equation*}
    \begin{vmatrix}
      1 & 0 & 0 & 2 \\
      0 & 3 & 4 & 5 \\
      5 & 4 & 0 & 3 \\
      2 & 0 & 0 & 1
    \end{vmatrix}
  \end{equation*}
\end{ex}



\begin{ex}
  \(\begin{vmatrix} 2x & x & 1 & 2 \\ 1 & x & 1 & -1\\ 3 & 2 & x & 1 \\ 1 & 1 & 1& x\end{vmatrix}\)
  中 \(x^4\), \(x^3\) 系数.
\end{ex}



\begin{ex}
  设方阵 \(A,B\) 满足 \(AB=BA\). 证明
  \(\det\begin{bmatrix}A & B\\-B & A\end{bmatrix}=\det(A^2+B^2)\).
\end{ex}


\begin{ex}
  \begin{enumerate}
  \item
    设 \(A\) 是正交矩阵, \(\det A < 0\). 证明 \(I_n + A\) 不可逆.
  \item   设 \(A\) 为奇数阶正交矩阵, \(\det A > 0\).
    证明 \(I_n - A\) 不可逆.
  \end{enumerate}
\end{ex}



\begin{ex}[\(\heartsuit\)]
  设 \(A\) 为对角占优方阵, 对角元素全为正数. 证明 \(\det A > 0\).
\end{ex}



\begin{ex}
  设 \(A\) 是元素为整数的可逆矩阵. 证明, 若 \(A\) 可逆, 并且 \(A^{-1}\)
  的元素也为整数的必要且充分条件是 \(|\det A| =  1\).
\end{ex}


\end{document}
